\chapter{Anomaly Scoring}
\label{ch:anomaly}

Path ranking is the second half of what the engine does: after the
matcher returns all paths satisfying a hypothesis, the
\Scorer{} assigns each path a composite score in $[0, 1]$ and
optional per-component breakdown so the analyst can triage. The
scorer is a weighted average of five dimensions
(entity rarity, edge rarity, neighbourhood concentration, temporal
novelty, and GNN threat), each computed as an average over the
vertices of the path.

\section{State and finalisation}

\begin{lstlisting}[style=rust]
pub struct AnomalyScorer {
    entity_freq:       HashMap<String, u64>,
    pair_freq:         HashMap<(String, String), u64>,
    first_seen:        HashMap<String, i64>,
    nc_cache:          HashMap<String, f64>,
    gnn_cache:         HashMap<String, f64>,
    log_max_freq:      f64,
    log_max_pair_freq: f64,
    tau_min:           i64,
    tau_max:           i64,
    weights:           ScoringWeights,
    finalized:         bool,
}
\end{lstlisting}

The scorer is populated incrementally during ingest
(\code{observe\_entity}, \code{observe\_edge}); those are
$\Ord{1}$ amortised and do not compute anything derived. After
ingest, \code{finalize} converts the raw counts into the
normalisation factors the per-path scoring needs.

\begin{algorithm}
\caption{\textsc{Anomaly-Finalize}}\label{alg:anomaly-finalize}
\KwIn{scorer $S$, graph $G$}
$S.\text{log\_max\_freq} \gets \ln(\max \text{entity\_freq.values})$\;
$S.\text{log\_max\_pair\_freq} \gets \ln(\max \text{pair\_freq.values})$\;
$(\tau_{\min}, \tau_{\max}) \gets (\min, \max)$ of the non-zero
values of \code{first\_seen.values}\;
$S.\text{nc\_cache} \gets \emptyset$\;
\ForEach{$v \in G.\text{entities}$}{
    $S.\text{nc\_cache}[v] \gets $ \textsc{Compute-NC}($v, G$)\;
}
$S.\text{finalized} \gets $ true\;
\end{algorithm}

\begin{complexity}
$\Ord{V + E + \sum_v d(v)} = \Ord{V + E}$ for the counters. The
Shannon-entropy cache costs an additional
$\Ord{\sum_v d(v)} = \Ord{E}$. Space is
$\Ord{V}$ for the caches plus whatever the raw counters already
consumed.
\end{complexity}

\section{Neighbourhood concentration}

\begin{algorithm}
\caption{\textsc{Compute-NC}}\label{alg:compute-nc}
\KwIn{vertex $v$, graph $G$}
\KwOut{neighbourhood concentration $\mathrm{NC}(v) \in [0, 1]$}
$\mathrm{counts} \gets \emptyset$; $\mathrm{total} \gets 0$\;
\ForEach{outgoing edge $(v, u) \in G$}{
    $\mathrm{counts}[\ell_V(u)] \mathrel{+}= 1$;
    $\mathrm{total} \mathrel{+}= 1$\;
}
\ForEach{incoming edge $(w, v) \in G$}{
    $\mathrm{counts}[\ell_V(w)] \mathrel{+}= 1$;
    $\mathrm{total} \mathrel{+}= 1$\;
}
$K \gets |\mathrm{counts}|$\;
\If{$K \leq 1$}{\Return 1.0\;}
$H \gets 0$\;
\ForEach{$c \in \mathrm{counts.values}$}{
    $p \gets c / \mathrm{total}$\;
    $H \gets H - p \log_2 p$\;
}
$H_{\max} \gets \log_2 K$; $\mathrm{NC} \gets 1 - H / H_{\max}$\;
\Return clamp($\mathrm{NC}$, 0, 1)\;
\end{algorithm}

$\mathrm{NC}$ is 1 when the neighbourhood is maximally homogeneous
(a single entity-type family) and 0 when maximally diverse; the
clamp guards floating-point drift near the extremes.

\begin{complexity}
$\Ord{d(v)}$ time, $\Ord{|\mathcal{T}_V|}$ space (neighbour
type counts only; small constant).
\end{complexity}

\section{Scoring a path}

Three of the four non-GNN components are simple averages over the
path's vertices; the edge-rarity component averages over the path's
edges.

\begin{algorithm}
\caption{\textsc{Score-Path}}\label{alg:score-path}
\KwIn{path $\pi = (v_0, \ldots, v_\ell)$, scorer $S$}
\KwOut{composite score $c$, breakdown $b$}
\If{$\pi$ empty or $\neg S.\text{finalized}$}{\Return $(0, \mathbf{0})$\;}
\tcp{Component 1: entity rarity}
$\mathrm{ER} \gets \frac{1}{\ell+1} \sum_v \left(1 - \frac{\ln(\text{entity\_freq}[v])}{S.\text{log\_max\_freq}}\right)$\;
\tcp{Component 2: edge rarity}
$\mathrm{EdgeR} \gets \frac{1}{\ell} \sum_{i=1}^{\ell} \left(1 - \frac{\ln(\text{pair\_freq}[v_{i-1}, v_i])}{S.\text{log\_max\_pair\_freq}}\right)$\;
\tcp{Component 3: NC average}
$\mathrm{NC} \gets \frac{1}{\ell+1} \sum_v S.\text{nc\_cache}[v]$\;
\tcp{Component 4: temporal novelty}
$\mathrm{TN} \gets \frac{1}{\ell+1} \sum_v \frac{\text{first\_seen}[v] - \tau_{\min}}{\tau_{\max} - \tau_{\min}}$\;
\tcp{Component 5: GNN threat}
$\mathrm{GT} \gets \frac{1}{\ell+1} \sum_v S.\text{gnn\_cache}.\text{get}(v).\text{unwrap\_or}(0)$\;
\tcp{composite via renormalised weights}
$W \gets \sum \text{weights}$; $c \gets \frac{1}{W} \sum_i w_i \cdot \text{component}_i$\;
\Return (clamp($c$, 0, 1), breakdown)\;
\end{algorithm}

Each component is individually clamped to $[0, 1]$ so that a stale
cache entry (extremely rare, but possible during a re-finalize
race) cannot make a component negative. The weight renormalisation
means the caller can set weights in any absolute scale.

\begin{complexity}
$\Ord{\ell}$ time, $\Ord{1}$ auxiliary space. The expensive state
was built by \textsc{Anomaly-Finalize}; per-path scoring is cheap.
\end{complexity}

\section{The trait seam (P2-A)}

The refactored scoring module
(\filepath{graph\_hunter\_core/src/scoring/}) makes each component a
separate \code{impl ScoreComponent}:

\begin{lstlisting}[style=rust]
pub trait ScoreComponent: Send + Sync {
    fn name(&self) -> &'static str;
    fn score_path(&self, ctx: &PathContext<'_>) -> f64;
    fn explain(&self, value: f64, ctx: &PathContext<'_>) -> Option<String> {
        None
    }
}
\end{lstlisting}

\code{CompositeScorer} holds a
\code{Vec<Box<dyn ScoreComponent>>} plus matching weights and
applies \textsc{Score-Path}-style aggregation. Adding a sixth
component (Bayesian, learned model, etc.) is a new struct + one
\code{vec!{}} push.

\begin{theorem}[v1/v2 equivalence]
For every weight vector $w$ and every path $\pi$,
\code{AnomalyScorer::score\_path}$(w, \pi)$ and
\code{AnomalyScorer::score\_path\_v2}$(w, \pi)$ agree to within
$8 \cdot \varepsilon_{\text{f64}}$.
\end{theorem}

This is verified by \code{tests/scoring\_snapshot.rs}. The
equivalence is the precondition for eventually retiring the hardcoded
v1 path.

\section{Node-level fast estimate}

During smart-DFS pruning (Chapter~\ref{ch:matcher}), we need an
$\Ord{1}$ estimate of a single vertex's anomaly so the running
\code{path\_anomaly\_sum} can reject branches. That's
\code{node\_anomaly\_estimate}:

\begin{lstlisting}[style=rust]
pub fn node_anomaly_estimate(&self, v: &str) -> f64 {
    if !self.finalized { return 0.5; }
    let er = /* Component 1 for a single vertex */;
    let nc = self.nc_cache.get(v).copied().unwrap_or(1.0);
    let tn = /* Component 4 for a single vertex */;
    let gt = self.gnn_cache.get(v).copied().unwrap_or(0.0);
    combine_with_weights(er, nc, tn, gt, &self.weights)
}
\end{lstlisting}

Edge rarity isn't part of the estimate because the DFS hasn't yet
chosen the next edge when it queries the bound. The estimate is
therefore a loose upper bound on the node's contribution to the
composite; it's admissible (never over-prunes) provided the weights
are non-negative, which the engine enforces at config load time.

\begin{inthecode}
\filepath{graph\_hunter\_core/src/anomaly.rs:148} ---
\code{finalize}.\\
\filepath{graph\_hunter\_core/src/anomaly.rs:191} ---
\code{compute\_nc}.\\
\filepath{graph\_hunter\_core/src/anomaly.rs:241} ---
\code{score\_path} (v1).\\
\filepath{graph\_hunter\_core/src/anomaly.rs:351} ---
\code{node\_anomaly\_estimate}.\\
\filepath{graph\_hunter\_core/src/scoring/} --- trait-based v2.\\
\filepath{graph\_hunter\_core/tests/scoring\_snapshot.rs} ---
v1/v2 equivalence pinning.
\end{inthecode}
